\section{Idea de negocio}

\subsection{Servicio ofrecido}

Pulso es un sistema en línea para organizar la información diaria de un
gimnasio. Permite registrar miembros, membresías, pagos, asistencia y planes de
entrenamiento. También permite consultar el progreso. Así se evita buscar la
misma información en cuadernos, hojas de cálculo y varias aplicaciones.

La versión inicial se concentra en los procesos que más valor aportan a la
muestra consultada: cobros, vigencias, asistencia, planes de entrenamiento y
progreso. Las pasarelas de pago y el acceso físico automatizado forman parte
del crecimiento posterior del producto. Pulso organiza información operativa y
deportiva; no sustituye un sistema contable ni un servicio médico.

\subsection{Cliente y personas usuarias}

El cliente principal es la persona dueña o administradora de un gimnasio. Esa
persona decide si contrata el servicio. Los entrenadores y los miembros usan el
sistema, pero no siempre son quienes lo compran. También se puede atender a un
entrenador que administra sus propios clientes.

Pulso puede atender gimnasios de distintos lugares porque funciona en línea.
Pérez Zeledón es el lugar de la validación inicial. La cercanía ayuda a realizar
la primera demostración y a acompañar el piloto, pero no limita el mercado.

La encuesta logró cinco respuestas del segmento responsable: tres de dueños o
administradores y dos de entrenadores con responsabilidades administrativas.
Dos administran establecimientos de 76 a 150 miembros, uno más de 150 y dos
indicaron no conocer el tamaño. Las herramientas informadas abarcan papel, hojas
de cálculo, WhatsApp, software especializado y sistemas internos. Esta variedad
aconseja ofrecer una incorporación flexible, pero la cantidad de casos no permite
estimar el tamaño del mercado.

\subsection{Forma de comercialización}

La venta comenzará con contacto directo y una demostración breve. Después
se ofrece un piloto de cuatro semanas por CRC 15\,000, con configuración y
acompañamiento incluidos. Al finalizar, el gimnasio puede continuar con el plan
que se ajuste a su uso y acompañamiento requerido. Se utilizarán una página informativa,
WhatsApp Business, correo y contenido demostrativo en redes sociales. Esta ruta
responde a la encuesta: dos responsables aceptaron probar la plataforma y cuatro
prefirieron conocerla antes de elegir un precio. La demostración y el piloto
permitirán conocer cómo se usa el sistema y qué mejoras necesitan los clientes.
Para iniciar esta validación en Pérez Zeledón, Pulso cuenta con Pit Bull Gym, en
Mollejones, como gimnasio aliado para la prueba inicial. Esta alianza facilita
acercar el sistema a un entorno real y obtener observaciones de una persona
administradora, lo que constituye una ventaja para la validación local.

\subsection{Necesidad atendida}

La propuesta busca ordenar clientes, cobros, vencimientos, asistencia y planes
de entrenamiento. No se trata solo de pasar un cuaderno a una pantalla. Cada
entrenador puede ver a sus miembros. Cada miembro puede ver su plan, su
membresía, sus pagos y su asistencia. En la encuesta, cuatro responsables
mencionaron problemas con cobros o pagos. Entre los miembros, seis señalaron los
vencimientos y cinco el progreso. Por eso la demostración inicia con esas
funciones.

\subsection{Impacto esperado}

Se espera mejorar la organización del gimnasio y la comunicación con sus
miembros. Durante el piloto se medirá el tiempo para registrar tareas, los
errores, las tareas completadas y la frecuencia de uso. Usar menos formularios
físicos también puede reducir el consumo de papel.

\section{Modelo de negocio}

Pulso propone una suscripción mensual según el uso y la ayuda que necesite el
gimnasio. El pago incluirá acceso a la plataforma, actualizaciones y soporte
inicial. El negocio se organiza en cuatro aspectos:

\begin{itemize}
  \item \textbf{Cliente}: responsables de gimnasios y entrenadores con necesidad
  de centralizar información.
  \item \textbf{Oferta}: administración y seguimiento deportivo en una sola
  plataforma adaptable al celular y la computadora.
  \item \textbf{Infraestructura}: equipo de desarrollo, aplicación, base de datos,
  servicios de nube, soporte y procesos de seguridad.
  \item \textbf{Viabilidad económica}: suscripciones recurrentes que deben cubrir
  infraestructura, soporte, mantenimiento e inversión de trabajo.
\end{itemize}

El Canvas resume a quién se ofrecerá el servicio, cómo se atenderá al cliente y
de dónde saldrán los ingresos para cubrir los gastos. Sus nueve bloques siguen
la estructura de ExpoTÉCNICA \parencite{mep2026} y se apoyan en el prototipo y
la encuesta. El piloto permitirá revisar si los gimnasios desean continuar
después de probar Pulso.

\begin{table}[htbp]
\caption{Resumen del modelo de negocio de Pulso}
\begin{tabularx}{\textwidth}{>{\bfseries}p{3.5cm}X}
\toprule
Elemento & Definición inicial \\
\midrule
Servicio & Sistema integral en línea de gestión administrativa y seguimiento deportivo. \\
Comprador & Dueño o administrador de gimnasio. \\
Usuarios & Responsables y entrenadores mediante el rol entrenador; miembros mediante su propio rol. \\
Ingreso & Escenario proyectado: un plan Starter de CRC 15\,000 y planes Pro Gym de CRC 25\,000 mensuales. \\
Escalamiento & Premium (CRC 45\,000) y Enterprise (desde CRC 70\,000) se consideran alternativas posteriores, sujetas a cotización y validación; no forman parte de la proyección financiera. \\
Entrega & Acceso web responsive con acompañamiento inicial. \\
Personalización & Nombre, logo y colores por gimnasio en la evolución comercial. \\
Mercado & Gimnasios; Pérez Zeledón es la validación inicial. \\
Madurez & Producto mínimo viable funcional, preparado para demostración y piloto. \\
\bottomrule
\end{tabularx}
\caption*{\fuentepropia}
\end{table}
